
%==========================================
%
% SIBGRAPI 2026 paper template
% Example of IEEEtran.cls
%
%==========================================

% *** Authors should verify (and, if needed, correct) their LaTeX system  ***
% *** with the testflow diagnostic prior to trusting their LaTeX platform ***
% *** with prodUCB1ion work. The IEEE's font choices and paper sizes can   ***
% *** trigger bugs that do not appear when using other class files.       ***                          ***
% The testflow support page is at:
% http://www.michaelshell.org/tex/testflow/

\documentclass[10pt,conference]{IEEEtran}


% Some very useful LaTeX packages include:
% (uncomment the ones you want to load)


% *** MISC UTILITY PACKAGES ***
%
%\usepackage{ifpdf}
% Heiko Oberdiek's ifpdf.sty is very useful if you need conditional
% compilation based on whether the output is pdf or dvi.
% usage:
% \ifpdf
%   % pdf code
% \else
%   % dvi code
% \fi
% The latest version of ifpdf.sty can be obtained from:
% http://www.ctan.org/pkg/ifpdf
% Also, note that IEEEtran.cls V1.7 and later provides a builtin
% \ifCLASSINFOpdf conditional that works the same way.
% When switching from latex to pdflatex and vice-versa, the compiler may
% have to be run twice to clear warning/error messages.






% *** CITATION PACKAGES ***
%
\usepackage{cite}
% cite.sty was written by Donald Arseneau
% V1.6 and later of IEEEtran pre-defines the format of the cite.sty package
% \cite{} output to follow that of the IEEE. Loading the cite package will
% result in citation numbers being automatically sorted and properly
% "compressed/ranged". e.g., [1], [9], [2], [7], [5], [6] without using
% cite.sty will become [1], [2], [5]--[7], [9] using cite.sty. cite.sty's
% \cite will automatically add leading space, if needed. Use cite.sty's
% noadjust option (cite.sty V3.8 and later) if you want to turn this off
% such as if a citation ever needs to be enclosed in parenthesis.
% cite.sty is already installed on most LaTeX systems. Be sure and use
% version 5.0 (2009-03-20) and later if using hyperref.sty.
% The latest version can be obtained at:
% http://www.ctan.org/pkg/cite
% The documentation is contained in the cite.sty file itself.






% *** GRAPHICS RELATED PACKAGES ***
%
\ifCLASSINFOpdf
   \usepackage[pdftex]{graphicx}
  % declare the path(s) where your graphic files are
   \graphicspath{{figs/}}
  % and their extensions so you won't have to specify these with
  % every instance of \includegraphics
   \DeclareGraphicsExtensions{.pdf,.jpeg,.png}
\else
  % or other class option (dvipsone, dvipdf, if not using dvips). graphicx
  % will default to the driver specified in the system graphics.cfg if no
  % driver is specified.
   \usepackage[dvips]{graphicx}
  % declare the path(s) where your graphic files are
   \graphicspath{{../figs/}}
  % and their extensions so you won't have to specify these with
  % every instance of \includegraphics
   \DeclareGraphicsExtensions{.eps}
\fi
% graphicx was written by David Carlisle and Sebastian Rahtz. It is
% required if you want graphics, photos, etc. graphicx.sty is already
% installed on most LaTeX systems. The latest version and documentation
% can be obtained at: 
% http://www.ctan.org/pkg/graphicx
% Another good source of documentation is "Using Imported Graphics in
% LaTeX2e" by Keith Reckdahl which can be found at:
% http://www.ctan.org/pkg/epslatex
%
% latex, and pdflatex in dvi mode, support graphics in encapsulated
% postscript (.eps) format. pdflatex in pdf mode supports graphics
% in .pdf, .jpeg, .png and .mps (metapost) formats. Users should ensure
% that all non-photo figures use a vector format (.eps, .pdf, .mps) and
% not a bitmapped formats (.jpeg, .png). The IEEE frowns on bitmapped formats
% which can result in "jaggedy"/blurry rendering of lines and letters as
% well as large increases in file sizes.
%
% You can find documentation about the pdfTeX application at:
% http://www.tug.org/applications/pdftex





% *** MATH PACKAGES ***
%
\usepackage[cmex10]{amsmath}
% A popular package from the American Mathematical Society that provides
% many useful and powerful commands for dealing with mathematics.
%
% Note that the amsmath package sets \interdisplaylinepenalty to 10000
% thus preventing page breaks from occurring within multiline equations. Use:
\interdisplaylinepenalty=2500
% after loading amsmath to restore such page breaks as IEEEtran.cls normally
% does. amsmath.sty is already installed on most LaTeX systems. The latest
% version and documentation can be obtained at:
% http://www.ctan.org/pkg/amsmath
\usepackage{amsthm}
\newtheorem{definition}{Definition}




% *** SPECIALIZED LIST PACKAGES ***
%
\usepackage{algorithmic}
% algorithmic.sty was written by Peter Williams and Rogerio Brito.
% This package provides an algorithmic environment fo describing algorithms.
% You can use the algorithmic environment in-text or within a figure
% environment to provide for a floating algorithm. Do NOT use the algorithm
% floating environment provided by algorithm.sty (by the same authors) or
% algorithm2e.sty (by Christophe Fiorio) as the IEEE does not use dedicated
% algorithm float types and packages that provide these will not provide
% correct IEEE style captions. The latest version and documentation of
% algorithmic.sty can be obtained at:
% http://www.ctan.org/pkg/algorithms
% Also of interest may be the (relatively newer and more customizable)
% algorithmicx.sty package by Szasz Janos:
% http://www.ctan.org/pkg/algorithmicx




% *** ALIGNMENT PACKAGES ***
%
\usepackage{array}
% Frank Mittelbach's and David Carlisle's array.sty patches and improves
% the standard LaTeX2e array and tabular environments to provide better
% appearance and additional user controls. As the default LaTeX2e table
% generation code is lacking to the point of almost being broken with
% respect to the quality of the end results, all users are strongly
% advised to use an enhanced (at the very least that provided by array.sty)
% set of table tools. array.sty is already installed on most systems. The
% latest version and documentation can be obtained at:
% http://www.ctan.org/pkg/array


% IEEEtran contains the IEEEeqnarray family of commands that can be used to
% generate multiline equations as well as matrices, tables, etc., of high
% quality.




% *** SUBFIGURE PACKAGES ***
\ifCLASSOPTIONcompsoc
  \usepackage[caption=false,font=normalsize,labelfont=sf,textfont=sf]{subfig}
\else
  \usepackage[caption=false,font=footnotesize]{subfig}
\fi
% subfig.sty, written by Steven Douglas Cochran, is the modern replacement
% for subfigure.sty, the latter of which is no longer maintained and is
% incompatible with some LaTeX packages including fixltx2e. However,
% subfig.sty requires and automatically loads Axel Sommerfeldt's caption.sty
% which will override IEEEtran.cls' handling of captions and this will result
% in non-IEEE style figure/table captions. To prevent this problem, be sure
% and invoke subfig.sty's "caption=false" package option (available since
% subfig.sty version 1.3, 2005/06/28) as this is will preserve IEEEtran.cls
% handling of captions.
% Note that the Computer Society format requires a larger sans serif font
% than the serif footnote size font used in traditional IEEE formatting
% and thus the need to invoke different subfig.sty package options depending
% on whether compsoc mode has been enabled.
%
% The latest version and documentation of subfig.sty can be obtained at:
% http://www.ctan.org/pkg/subfig

\usepackage[utf8]{inputenc}
\usepackage{listingsutf8}

% *** FLOAT PACKAGES ***
%
%\usepackage{fixltx2e}
% fixltx2e, the successor to the earlier fix2col.sty, was written by
% Frank Mittelbach and David Carlisle. This package corrects a few problems
% in the LaTeX2e kernel, the most notable of which is that in current
% LaTeX2e releases, the ordering of single and double column floats is not
% guaranteed to be preserved. Thus, an unpatched LaTeX2e can allow a
% single column figure to be placed prior to an earlier double column
% figure.
% Be aware that LaTeX2e kernels dated 2015 and later have fixltx2e.sty's
% corrections already built into the system in which case a warning will
% be issued if an attempt is made to load fixltx2e.sty as it is no longer
% needed.
% The latest version and documentation can be found at:
% http://www.ctan.org/pkg/fixltx2e


%\usepackage{stfloats}
% stfloats.sty was written by Sigitas Tolusis. This package gives LaTeX2e
% the ability to do double column floats at the bottom of the page as well
% as the top. (e.g., "\begin{figure*}[!b]" is not normally possible in
% LaTeX2e). It also provides a command:
%\fnbelowfloat
% to enable the placement of footnotes below bottom floats (the standard
% LaTeX2e kernel puts them above bottom floats). This is an invasive package
% which rewrites many portions of the LaTeX2e float routines. It may not work
% with other packages that modify the LaTeX2e float routines. The latest
% version and documentation can be obtained at:
% http://www.ctan.org/pkg/stfloats
% Do not use the stfloats baselinefloat ability as the IEEE does not allow
% \baselineskip to stretch. Authors submitting work to the IEEE should note
% that the IEEE rarely uses double column equations and that authors should try
% to avoid such use. Do not be tempted to use the cuted.sty or midfloat.sty
% packages (also by Sigitas Tolusis) as the IEEE does not format its papers in
% such ways.
% Do not attempt to use stfloats with fixltx2e as they are incompatible.
% Instead, use Morten Hogholm'a dblfloatfix which combines the features
% of both fixltx2e and stfloats:
%
% \usepackage{dblfloatfix}
% The latest version can be found at:
% http://www.ctan.org/pkg/dblfloatfix
\usepackage{float}



\usepackage{tikz}
\usepackage{tabularx}
\usepackage{booktabs}


% *** PDF, URL AND HYPERLINK PACKAGES ***
%
\usepackage{url}
% url.sty was written by Donald Arseneau. It provides better support for
% handling and breaking URLs. url.sty is already installed on most LaTeX
% systems. The latest version and documentation can be obtained at:
% http://www.ctan.org/pkg/url
% Basically, \url{my_url_here}.




% *** Do not adjust lengths that control margins, column widths, etc. ***
% *** Do not use packages that alter fonts (such as pslatex).         ***
% There should be no need to do such things with IEEEtran.cls V1.6 and later.
% (Unless specifically asked to do so by the journal or conference you plan
% to submit to, of course. )


% correct bad hyphenation here
\hyphenation{op-tical net-works semi-conduc-tor}

\lstset{
    inputencoding=utf8,
    extendedchars=true,
    literate=
      % Minúsculas acentuadas
      {á}{{\'a}}1 {é}{{\'e}}1 {í}{{\'i}}1 {ó}{{\'o}}1 {ú}{{\'u}}1
      {à}{{\`a}}1 {è}{{\`e}}1 {ì}{{\`i}}1 {ò}{{\`o}}1 {ù}{{\`u}}1
      {â}{{\^a}}1 {ê}{{\^e}}1 {î}{{\^i}}1 {ô}{{\^o}}1 {û}{{\^u}}1
      {ã}{{\~a}}1 {õ}{{\~o}}1
      {ç}{{\c{c}}}1
      % Maiúsculas acentuadas
      {Á}{{\'A}}1 {É}{{\'E}}1 {Í}{{\'I}}1 {Ó}{{\'O}}1 {Ú}{{\'U}}1
      {À}{{\`A}}1 {È}{{\`E}}1 {Ì}{{\`I}}1 {Ò}{{\`O}}1 {Ù}{{\`U}}1
      {Â}{{\^A}}1 {Ê}{{\^E}}1 {Î}{{\^I}}1 {Ô}{{\^O}}1 {Û}{{\^U}}1
      {Ã}{{\~A}}1 {Õ}{{\~O}}1
      {Ç}{{\c{C}}}1
}

%------------------------------------------------------------------------- 
% change the % on next lines to produce the final camera-ready version 
\newif\iffinal
%\finalfalse
\finaltrue
\newcommand{\cmtid}{99999}
%------------------------------------------------------------------------- 

\iffinal
\else
\usepackage[switch]{lineno}
\fi

\begin{document}
%
% paper title
% Titles are generally capitalized except for words such as a, an, and, as,
% at, but, by, for, in, nor, of, on, or, the, to and up, which are usually
% not capitalized unless they are the first or last word of the title.
% Linebreaks \\ can be used within to get better formatting as desired.
% Do not put math or special symbols in the title.
\title{Modelagem Ontológica e Processamento Semântico para Ecologia de Estradas no Banhado do Taim}


% author names and affiliations
% use a multiple column layout for up to two different
% affiliations

\iffinal

% author names and affiliations
% use a multiple column layout for up to three different
% affiliations
\author{\IEEEauthorblockN{Guilherme M. Einloft and Gabriel Stiegemeier}
\IEEEauthorblockA{Ciência da Computação\\
Universidade Federal de Santa Maria\\
Santa Maria, RS, Brazil 97105-900\\
Email: gmeinloft@inf.ufsm.br and gstiegemeier@inf.ufsm.br}}

% conference papers do not typically use \thanks and this command
% is locked out in conference mode. If really needed, such as for
% the acknowledgment of grants, issue a \IEEEoverridecommandlockouts
% after \documentclass

% for over three affiliations, or if they all won't fit within the width
% of the page, use this alternative format:
% 
%\author{\IEEEauthorblockN{Michael Shell\IEEEauthorrefmark{1},
%Homer Simpson\IEEEauthorrefmark{2},
%James Kirk\IEEEauthorrefmark{3}, 
%Montgomery Scott\IEEEauthorrefmark{3} and
%Eldon Tyrell\IEEEauthorrefmark{4}}
%\IEEEauthorblockA{\IEEEauthorrefmark{1}School of Electrical and Computer Engineering\\
%Georgia Institute of Technology,
%Atlanta, Georgia 30332--0250\\ Email: see http://www.michaelshell.org/contact.html}
%\IEEEauthorblockA{\IEEEauthorrefmark{2}Twentieth Century Fox, Springfield, USA\\
%Email: homer@thesimpsons.com}
%\IEEEauthorblockA{\IEEEauthorrefmark{3}Starfleet Academy, San Francisco, California 96678-2391\\
%Telephone: (800) 555--1212, Fax: (888) 555--1212}
%\IEEEauthorblockA{\IEEEauthorrefmark{4}Tyrell Inc., 123 Replicant Street, Los Angeles, California 90210--4321}}

\else
  \author{SIBGRAPI Paper ID: \cmtid \\ }
  \linenumbers
\fi


% make the title area
\maketitle

% As a general rule, do not put math, special symbols or citations
% in the abstract


% no keywords




% For peerreview papers, this IEEEtran command inserts a page break and
% creates the second title. It will be ignored for other modes.
\IEEEpeerreviewmaketitle

\begin{abstract}
Este relatório descreve o desenvolvimento de uma base de conhecimento ontológica voltada para o domínio ecológico e operacional do Banhado do Taim, com foco na redução do atropelamento de fauna silvestre. A arquitetura foi construída no padrão formal OWL 2, estabelecendo uma representação focada nos eventos de atropelamento, nos animais afetados e nas características das rodovias locais. O projeto integrou a extração automática de informações a partir de múltiplas fontes de dados (como Wikipedia, Wikidata e IBGE), utilizando técnicas de Processamento de Linguagem Natural (PLN) para limpar, filtrar e transformar textos dispersos em relações estruturadas. Adicionalmente, desenvolveu-se um módulo de automação para o povoamento sintético da ontologia, alcançando um bom volume de instâncias interconectadas. O modelo gerado é validado por um conjunto de 30 consultas operacionais SPARQL e propõe-se uma base conceitual para a futura integração desse modelo com sistemas de Inteligência Artificial Explicável (XAI).

\end{abstract}

\section{Introdução e Contextualização do Problema}

O atropelamento de fauna silvestre em estradas que cortam ecossistemas sensíveis representa um impacto humano negativo à biodiversidade local. A Estação Ecológica (ESEC) do Taim, localizada no bioma Pampa e recortada pela rodovia BR-471, configura-se como um cenário de vulnerabilidade ecológica. O tráfego interurbano dividindo o ecossistema causa a mortalidade de vertebrados, o efeito de barreira e a separação de habitats.

Para diminuir esses problemas, é necessário adotar métodos de gestão baseados em dados que consigam prever riscos. Contudo, a variação e a dispersão dos dados ecológicos, espaciais e de trânsito dificultam as análises amplas em bancos de dados tradicionais. Nesse contexto, este projeto propõe o desenvolvimento de um \textit{Knowledge Graph} (Grafo de Conhecimento) baseado em uma ontologia de domínio (OWL 2), capaz de organizar as interações entre os habitats, os animais, o tráfego e o clima. O objetivo principal é criar um sistema computacional que ajude a formular medidas de proteção localizadas, baseadas em cenários de previsão.

\section{Metodologia de Modelagem}

A criação da ontologia seguiu os princípios recomendados pelo método NeOn, priorizando o reaproveitamento de vocabulários, o alinhamento lógico e a clareza nas descrições. O esquema da ontologia (\textit{TBox}) foi codificado utilizando a biblioteca Python \texttt{owlready2}. 

A estrutura resultou em uma classificação organizada composta por mais de 20 classes principais (incluindo \texttt{Animal}, \texttt{Rodovia}, \texttt{EventoAtropelamento}, \texttt{CondicaoAmbiental}, \texttt{CaracteristicaHabitat} e \texttt{FatorRisco}), interligadas por 13 propriedades de objeto (\textit{Object Properties}) e 15 propriedades de dados (\textit{Data Properties}). Impôs-se rigor formal via declarações de regras matemáticas (como restrições existenciais e universais), separação clara entre grupos de classes e definição rigorosa dos pontos de partida e chegada (\textit{domain}/\textit{range}) para as relações. Duas categorias básicas, de ordem espacial (\texttt{localizadoEm}) e temporal (\texttt{ocorreDurantePeriodo}), foram modeladas explicitamente para garantir o suporte de tempo e espaço nas análises.

\begin{figure}[htbp]
\centering
\includegraphics[width=\columnwidth]{taim_classes.png}
\caption{Hierarquia de classes da ontologia.}
\label{fig:taim_classes}
\end{figure}

\subsection{Justificativas Teóricas para a Arquitetura}

A decisão fundamental na modelagem foi a adoção de uma arquitetura baseada em eventos. A classe \texttt{EventoAtropelamento} foi criada como o ponto central da rede, conectando os itens parados (ex: habitat, rodovia) e os em movimento (ex: animal, veículos) como parte de um único acontecimento. Esse formato de estrela melhora o desempenho em consultas com muitas variáveis ao mesmo tempo (ex: cruzar tipo de animal, velocidade da via e temperatura local). As relações funcionais (\texttt{FunctionalProperty}) foram muito usadas para impedir que conexões incorretas ocorram e para reforçar a lógica do modelo, garantindo limite de um-para-um em atributos rígidos como \texttt{ocorreEmRodovia}.

\subsection{Alternativas Arquiteturais Descartadas}

Durante as fases de elicitação, duas abordagens alternativas foram metodologicamente ponderadas e descartadas mediante critérios de custo computacional e adequabilidade semântica:

\begin{enumerate}
    \item \textbf{Abordagem 4D-Fluent:} O padrão \textit{4D-Fluent} sugere a divisão dos objetos em "fatias de tempo" para lidar com características que mudam. Isso foi rejeitado pelo aumento no número de dados gerados e pelo maior tempo nas buscas do sistema. Como a rodovia e a espécie não mudam de sentido ao longo do tempo no escopo de um atropelamento, a abordagem focada no evento (onde o próprio evento guarda a informação de tempo) mostrou-se mais simples e eficiente.
    
    \item \textbf{Reutilização da Ontologia SSN/SOSA:} A ideia inicial era usar a ontologia \textit{SOSA} para lidar com os dados de clima. No entanto, o excesso de detalhes dessa ontologia (que obriga a criação de sensores, atuadores e procedimentos) criava uma complexidade desnecessária para o nosso projeto prático. Optou-se pela criação direta da classe \texttt{CondicaoClimatica} ligada ao evento, o que diminui o tamanho do modelo sem perder informações importantes.
\end{enumerate}

\section{Extração e Povoamento da Ontologia}

\subsection{Fontes de Informação}

O planejamento inicial buscou várias fontes de informação para o povoamento da ontologia, categorizadas como: bases estruturadas (Wikidata, DBpedia, IBGE), fontes textuais não estruturadas (Wikipedia, artigos científicos, relatórios ambientais), dados geográficos (OpenStreetMap) e Modelos de Linguagem (LLMs). Aplicou-se uma etapa de testes de viabilidade sobre cada fonte candidata, avaliando-se a \textbf{acessibilidade técnica} (disponibilidade de APIs, ausência de \textit{paywalls}, conformidade de \textit{scraping}) e o \textbf{rendimento semântico} (volume de entidades e relações extraíveis para o domínio de atropelamento de fauna e ecologia do Taim).

As fontes que foram \textbf{aprovadas} e efetivamente integradas ao \textit{pipeline} de extração são:

\begin{enumerate}[label=\textbf{\arabic*.}]
    \item \textbf{Wikipedia} (fontes textuais não estruturadas): Páginas enciclopédicas sobre a Estação Ecológica do Taim, Lagoa Mirim, rodovia BR-471, zonas úmidas (banhados) e fichas de espécies (\textit{Cerdocyon thous}, \textit{Caiman latirostris}, \textit{Hydrochoerus hydrochaeris}). Foi a fonte responsável pela maioria das triplas extraídas via PLN.

    \item \textbf{Wikidata} (base estruturada via SPARQL): Consultas ao \textit{endpoint} SPARQL do Wikidata para obtenção de dados estruturados sobre espécies faunísticas relevantes — nome científico, nome vernáculo, status de conservação IUCN e habitat. Os dados foram convertidos em parágrafos textuais e integrados ao \textit{pipeline} NLP unificado.

    \item \textbf{IBGE — API de Localidades} (base governamental estruturada): Consultas à API REST do IBGE para extração de dados geográfico-administrativos dos municípios de Rio Grande e Santa Vitória do Palmar (RS), incluindo microrregião, mesorregião e região imediata. Estes dados contextuais foram textualizados e processados para enriquecimento das entidades de localização geográfica.

    \item \textbf{Portais de Notícias Ambientais} (fontes textuais não estruturadas): Reportagens do portal G1 (Globo) e do jornalísmo ambiental O Eco, que apresentaram conteúdo textual rico sobre a ESEC Taim, políticas de conservação, eventos de atropelamento e medidas de mitigação.
\end{enumerate}

As fontes \textbf{rejeitadas} na triagem, com suas respectivas justificativas:

\begin{itemize}
    \item \textbf{Remoção por Inacessibilidade:} A DBpedia apresentou indisponibilidade persistente (erro HTTP 503 — \textit{Service Unavailable}) em todas as consultas SPARQL. O OpenStreetMap (Overpass API) retornou erro HTTP 406 (\textit{Not Acceptable}), bloqueando as rotinas automatizadas. Repositórios acadêmicos como o SciELO e o LUME/UFRGS retornaram erros 404 (\textit{Not Found}) para os artigos prospectados. LLMs (Modelos de Linguagem) foram descartados por ausência de API pública gratuita configurada no ambiente de execução.

    \item \textbf{Remoção por Baixo Rendimento Semântico:} O periódico universitário da UFPEL, embora acessível (HTTP 200), apresentou rendimento semântico insuficiente: durante a extração de teste, gerou apenas 1 termo de domínio e zero triplas ontológicas relevantes, inviabilizando sua inclusão no \textit{pipeline} frente ao custo computacional.
\end{itemize}

\subsection{Protocolos de Raspagem e Processamento de Linguagem Natural (PLN)}

A fase de ingestão empregou a biblioteca \texttt{requests} aliada ao \texttt{BeautifulSoup} para parsing de HTML. Como a obtenção de dados de portais públicos e enciclopédias requer práticas responsáveis, o sistema incluiu o tratamento automático de erros de conexão (como \texttt{HTTPError} e \texttt{Timeout}) e respeitou intervalos amigáveis entre os acessos (\textit{politeness delay}) para não sobrecarregar os servidores alvo.

\subsubsection{Reconhecimento de Entidades e Refinamento}

Em seguida, o processamento dos textos foi feito usando a biblioteca \texttt{spaCy} e o modelo linguístico pré-treinado \texttt{pt\_core\_news\_sm}. Devido à falta de palavras específicas de ecologia neste modelo geral, uma solução mista foi implementada. O modelo foi expandido com um componente \texttt{EntityRuler}, contendo um vocabulário controlado com os nomes taxonômicos da região (ex: \textit{Caiman latirostris}, \textit{Cerdocyon thous}) e os jargões locais do Banhado do Taim.

Para garantir que a extração resultasse em nomes de entidades limpos e padronizados, a árvore sintática (\textit{dependency tree}) de cada frase foi analisada. O algoritmo foi programado para varrer os nós das palavras e ativamente ignorar classes morfológicas indesejadas, como determinantes (artigos "o", "a", "um") e adposições (preposições "de", "em", "para"). Adicionalmente, aplicou-se uma limpeza por Expressões Regulares (RegEx) para remover qualquer sufixo ou prefixo residual, assegurando que o nome final da entidade estivesse gramaticalmente isolado (e.g., convertendo "da Lagoa Mirim" para "Lagoa Mirim").

\subsection{Regras de Extração e Instanciação}

O \textit{pipeline} consolidou a extração de triplas (Sujeito, Relação, Objeto) mediante duas vias complementares:
\begin{itemize}
    \item \textbf{Parsing de Dependência Sintática (\textit{Dependency Parsing}):} Reconhecimento dos padrões sintáticos (ex: sujeito nominal $\leftarrow$ verbo $\rightarrow$ objeto direto) e mapeamento explícito dos verbos para as propriedades ontológicas, como as propriedades \texttt{atravessaHabitat} e \texttt{possuiHabitat}.
    \item \textbf{Co-ocorrência Contextual:} Dedução heurística de relações baseada na adjacência sentencial de entidades de classes predefinidas, obedecendo às restrições de \textit{domain} e \textit{range} impostas pela \textit{TBox}.
\end{itemize}

Durante a construção do \textit{pipeline}, foi necessário aplicar regras de refinamento de linguagem natural: exigiu-se a correspondência correta de classe ontológica em ambas as extremidades das relações (sujeito e objeto) para evitar a formação de pares sem sentido. Além disso, o motor de leitura (\textit{tokens}) foi programado para rejeitar palavras irrelevantes (como artigos e preposições) e termos de domínio genéricos (e.g., ``trecho'', ``morte''), garantindo a extração de entidades nominais úteis.

Este artefato de extração deduplicado gerou 27 triplas, derivadas de 711 menções a entidades processadas e persistidas nos formatos JSON e CSV. A Tabela \ref{tab:triplas} ilustra um subconjunto destas triplas, demonstrando a precisão do mapeamento das entidades nativas para as classes ontológicas.

\begin{table*}[h]
\centering
\caption{Exemplos de triplas semânticas extraídas com sucesso pelo pipeline de PLN.}
\label{tab:triplas}
\begin{tabular}{@{}llll@{}}
\toprule
\textbf{Entidade Origem} & \textbf{Relação Ontológica} & \textbf{Entidade Destino} & \textbf{Classes Mapeadas} \\ \midrule
capivara & \texttt{possuiHabitat} & banhado & \texttt{Mamifero $\rightarrow$ TipoVegetacao} \\
atropelamentos & \texttt{envolveAnimal} & aves & \texttt{EventoAtropelamento $\rightarrow$ Ave} \\
br 471 & \texttt{atravessaHabitat} & banhado & \texttt{Rodovia $\rightarrow$ TipoVegetacao} \\
Lagoa Mirim & \texttt{ligado} & Lagoa Patos & \texttt{CorpoHidrico $\rightarrow$ CorpoHidrico} \\ 
jacaré & \texttt{possuiHabitat} & banhado & \texttt{Reptil $\rightarrow$ TipoVegetacao} \\ \bottomrule
\end{tabular}
\end{table*}

\subsection{Povoamento Automatizado da Ontologia}

Visando atingir a massa crítica e a expressividade necessárias para as consultas complexas, o requisito de instanciação de indivíduos (população da \textit{ABox}) foi preenchido através do desenvolvimento de um script automatizado em Python (\texttt{owlready2}). 

Esse script encarregou-se de gerar e relacionar 125 instâncias adicionais, simulando \texttt{EventosAtropelamento}, \texttt{Mamiferos}, \texttt{Aves}, \texttt{Répteis}, \texttt{TrechosRodoviarios} e \texttt{TiposVegetacao}. Ao final do processo, unindo o conhecimento básico extraído e a rotina de população sintética, a base (\texttt{ontologia\_taim\_povoada.owl}) totalizou 152 indivíduos devidamente interligados pelas relações da \textit{TBox} (como \texttt{envolveAnimal}, \texttt{ocorreEmTrecho} e \texttt{possuiHabitat}), formando a base de dados pronta para as etapas de análise seguintes.

A Tabela \ref{tab:distribuicao} quantifica a distribuição dos novos indivíduos criados automaticamente pelo script Python, garantindo o volume necessário de instâncias (\textit{ABox}) para simular um ambiente de banco de dados robusto.

\begin{table}[h]
\centering
\caption{Distribuição dos indivíduos sintéticos gerados pelo script de população.}
\label{tab:distribuicao}
\begin{tabular}{@{}lc@{}}
\toprule
\textbf{Classe Ontológica} & \textbf{Quantidade de Instâncias Geradas} \\ \midrule
\texttt{EventoAtropelamento} & 40 \\
\texttt{Mamifero} & 30 \\
\texttt{TrechoRodoviario} & 15 \\
\texttt{Reptil} & 15 \\
\texttt{Ave} & 15 \\
\texttt{TipoVegetacao} & 10 \\ \midrule
\textbf{Total Gerado} & \textbf{125} \\ \bottomrule
\end{tabular}
\end{table}

\section{Execução e Validação Semântica}

A validação do conhecimento inferido foi conduzida sistematicamente sobre a infraestrutura da linguagem \textit{SPARQL Protocol and RDF Query Language}. Formulou-se um corpus composto por 30 consultas, segmentadas rigidamente em cinco escopos funcionais, executadas e transacionadas via motor \texttt{rdflib}:

\begin{enumerate}
    \item \textbf{Consultas Simples (Filtragem por Classe):} Inventário matricial do domínio, e.g., quantificação das espécies avaliadas quanto aos status IUCN, atributos estáticos de rodovias e taxonomia animal subjacente.
    \item \textbf{Consultas Complexas (Grafo Polidimensional):} Buscas que atravessam vários relacionamentos ao mesmo tempo. Destacou-se o cruzamento de cadeias completas relacionando a Espécie $\rightarrow$ Animal $\rightarrow$ Evento $\rightarrow$ TrechoRodoviário $\rightarrow$ Habitat $\rightarrow$ Coordenadas Geográficas.
    \item \textbf{Filtros Operacionais:} Introdução de delimitadores condicionais. Exemplificam-se os mapeamentos noturnos e a convergência de variáveis (pista $< 7.5$m convergindo com alta volumetria veicular).
    \item \textbf{Consultas Analíticas de Agregação:} Operações envolvendo \texttt{GROUP BY}, \texttt{MAX}, \texttt{AVG} e \texttt{COUNT}. Essenciais para calcular as estatísticas gerais do modelo, como descobrir a diversidade de habitats das espécies ou identificar a faixa de temperatura em que mais ocorrem atropelamentos de répteis (e.g. 20-30$^\circ$C).
    \item \textbf{Cenários Operacionais Reais:} Resoluções aplicadas a demandas decisórias de gestores ambientais. Mapeou-se a prioridade prática para a instalação de estruturas de proteção (passagens inferiores), identificando locais onde existem fatores de risco mas não há nenhuma \texttt{MedidaMitigacao} cadastrada (\texttt{FILTER NOT EXISTS}).
\end{enumerate}

O comportamento empírico do protocolo atestou a ausência de ambiguidades formais, provando que o modelo ontológico é aderente à lógica de Primeira Ordem (\textit{FOL - First-Order Logic}) contida em seus metadados semânticos.

\section{Integração Preditiva e Explicabilidade}

Embora a ontologia funcione principalmente como um banco de dados organizado de forma inteligente, a sua estrutura permite a criação de soluções inovadoras usando inteligência artificial.

O modelo de previsão pode ser construído integrando a base \textit{OWL/SPARQL} com ferramentas de \textit{Machine Learning} probabilístico ou abordagens baseadas em regras (\textit{SWRL - Semantic Web Rule Language}). Utilizando o histórico de dados (\textit{features} como km da via, período sazonal, umidade e largura de pista), as redes neurais podem estimar um coeficiente de risco de novos atropelamentos nos "Hotspots" identificados.

Uma vantagem dessa abordagem, no entanto, está na capacidade de explicar as próprias conclusões (XAI - \textit{eXplainable Artificial Intelligence}). Enquanto modelos estatísticos agem como caixas-pretas (\textit{Black-Box}), a base ontológica permite a geração automática de justificativas, explicando as causas a partir dos dados estruturados. O sistema não apenas prevê de forma matemática onde ocorrerá um atropelamento; ele explica essa previsão mostrando como se cruzam a rota de migração do \textit{Cerdocyon thous} (fator ecológico), o habitat de área alagada (\textit{CorpoHidrico}) e a falta de barreiras de proteção no tráfego local (fatores humanos e métricas da rodovia).

\section{Conclusões}
O sistema computacional desenvolvido atende a todos os requisitos propostos, comprovando que o uso de tecnologias de Web Semântica é viável para organizar cenários ecológicos e problemas reais do meio ambiente. Através da construção de uma arquitetura ontológica no padrão OWL 2, tornou-se possível mapear as interações entre a fauna silvestre que habita a região do Banhado do Taim, os atributos espaciais da rodovia BR-471 e os vários fatores de risco operacionais que causam os atropelamentos. O trabalho superou a fase puramente teórica ao implementar soluções de código para alimentar a base ativamente.

O uso combinado de rotinas de coleta de dados (Web Scraping) e Processamento de Linguagem Natural (PLN) mostrou ser uma ferramenta poderosa para extrair informações soltas na internet e convertê-las de modo automático em um Grafo de Conhecimento útil e estruturado. A execução adicional de um script focado no povoamento rápido de novos indivíduos garantiu que o ambiente final não ficasse vazio, permitindo que a infraestrutura lógica do projeto fosse comprovadamente validada na prática por meio das consultas SPARQL. 

Como resultado de todas essas fases, modelagem, busca, refinamento gramatical e instânciação automatizada, entregamos um produto computacional sólido. A estrutura desse banco de dados pavimenta o caminho para a união com sistemas modernos de aprendizado de máquina, promovendo uma inteligência artificial capaz de prever ameaças e explicar cientificamente as suas próprias previsões com base nos relacionamentos do habitat.





% An example of a floating figure using the graphicx package.
% Note that \label must occur AFTER (or within) \caption.
% For figures, \caption should occur after the \includegraphics.
% Note that IEEEtran v1.7 and later has special internal code that
% is designed to preserve the operation of \label within \caption
% even when the captionsoff option is in effect. However, because
% of issues like this, it may be the safest practice to put all your
% \label just after \caption rather than within \caption{}.
%
% Reminder: the "draftcls" or "draftclsnofoot", not "draft", class
% option should be used if it is desired that the figures are to be
% displayed while in draft mode.
%
% Note that the IEEE typically puts floats only at the top, even when this
% results in a large percentage of a column being occupied by floats.


% An example of a double column floating figure using two subfigures.
% (The subfig.sty package must be loaded for this to work.)
% The subfigure \label commands are set within each subfloat command,
% and the \label for the overall figure must come after \caption.
% \hfil is used as a separator to get equal spacing.
% Watch out that the combined width of all the subfigures on a 
% line do not exceed the text width or a line break will occur.
%
%\begin{figure*}[!t]
%\centering
%\subfloat[Case I]{\includegraphics[width=2.4in]{figs/sibgrapi-2026.png}%
%\label{fig_first_case}}
%\hfil
%\subfloat[Case II]{\includegraphics[width=2.4in]{figs/sibgrapi-2026.png}%
%\label{fig_second_case}}
%\caption{SIBGRAPI - Conference on Graphics, Patterns and Images.}
%\label{fig_sim2}
%\end{figure*}
%
% Note that often IEEE papers with subfigures do not employ subfigure
% captions (using the optional argument to \subfloat[]), but instead will
% reference/describe all of them (a), (b), etc., within the main caption.
% Be aware that for subfig.sty to generate the (a), (b), etc., subfigure
% labels, the optional argument to \subfloat must be present. If a
% subcaption is not desired, just leave its contents blank,
% e.g., \subfloat[].


% An example of a floating table. Note that, for IEEE style tables, the
% \caption command should come BEFORE the table and, given that table
% captions serve much like titles, are usually capitalized except for words
% such as a, an, and, as, at, but, by, for, in, nor, of, on, or, the, to
% and up, which are usually not capitalized unless they are the first or
% last word of the caption. Table text will default to \footnotesize as
% the IEEE normally uses this smaller font for tables.
% The \label must come after \caption as always.
%
%\begin{table}[]
%% increase table row spacing, adjust to taste
%\renewcommand{\arraystretch}{1.3}
% if using array.sty, it might be a good idea to tweak the value of
% \extrarowheight as needed to properly center the text within the cells
%\caption{An Example of a Table}
%\label{table_example}
%\centering
%% Some packages, such as MDW tools, offer better commands for making tables
%% than the plain LaTeX2e tabular which is used here.
%\begin{tabular}{|c||c|}
%\hline
%One & Two\\
%\hline
%Three & Four\\
%\hline
%\end{tabular}
%\end{table}


% Note that the IEEE does not put floats in the very first column
% - or typically anywhere on the first page for that matter. Also,
% in-text middle ("here") positioning is typically not used, but it
% is allowed and encouraged for Computer Society conferences (but
% not Computer Society journals). Most IEEE journals/conferences use
% top floats exclusively. 
% Note that, LaTeX2e, unlike IEEE journals/conferences, places
% footnotes above bottom floats. This can be corrected via the
% \fnbelowfloat command of the stfloats package.



% trigger a \newpage just before the given reference
% number - used to balance the columns on the last page
% adjust value as needed - may need to be readjusted if
% the document is modified later
%\IEEEtriggeratref{8}
% The "triggered" command can be changed if desired:
%\IEEEtriggercmd{\enlargethispage{-5in}}

% references section

% can use a bibliography generated by BibTeX as a .bbl file
% BibTeX documentation can be easily obtained at:
% http://mirror.ctan.org/biblio/bibtex/contrib/doc/
% The IEEEtran BibTeX style support page is at:
% http://www.michaelshell.org/tex/ieeetran/bibtex/
%\bibliographystyle{IEEEtran}
% argument is your BibTeX string definitions and bibliography database(s)
%\bibliography{example}
%
% <OR> manually copy in the resultant .bbl file
% set second argument of \begin to the number of references
% (used to reserve space for the reference number labels box)
%\begin{thebibliography}{1}
%
%\bibitem{IEEEhowto:kopka}
%H.~Kopka and P.~W. Daly, \emph{A Guide to \LaTeX}, 3rd~ed.\hskip 1em plus
%  0.5em minus 0.4em\relax Harlow, England: Addison-Wesley, 1999.

%\end{thebibliography}




% that's all folks
\end{document}


